% from main.tex
%%%%%%%%%%  scan idx 242  |  folio 229  %%%%%%%%%%
pleine de \(D(X)\) formée des objets pseudo-cohérents relativement à \(f\).

On dira que \(f\) est \(n\)-pseudo-cohérent (resp. pseudo-cohérent) si \(\mathcal O_X\) est \(n\)-pseudo-cohérent (resp. pseudo-cohérent) relativement à \(f\). On dira parfois, quand cela ne pourra prêter à confusion, « \(n\)-pseudo-cohérent (resp. pseudo-cohérent) relativement à \(Y\) » au lieu de « relativement à \(f\) », et l'on écrira aussi \(D(X)_{Y\text{-}n\text{-}\coh}\) (resp. \(D(X)_{Y\text{-}\coh}\)) au lieu de \(D(f)_{n\text{-}\coh}\) (resp. \(D(f)_{\mathrm{coh}}\)). Enfin pour \(*=+,-,b\), on posera
\[
\begin{gathered}
  D^*(f)_{n\text{-}\coh}=D^*(X)\cap D(f)_{n\text{-}\coh},\\
  D^*(f)_{\mathrm{coh}}=D^*(X)\cap D(f)_{\mathrm{coh}}.
\end{gathered}
\]

\textbf{Remarques 1.2.1.}
\begin{enumerate}[label=\alph*)]
\item Les catégories \(D(f)_{n\text{-}\coh}\), \(D(f)_{\mathrm{coh}}\) ne sont pas en général contenues dans \(D^{-}(X)\). Il en est ainsi cependant si \(X\) est quasi-compact, car un objet de \(D(X)\) à cohomologie localement bornée supérieurement est alors à cohomologie bornée supérieurement.

\item Les objets de \(D(f)_{\mathrm{coh}}\) sont à cohomologie quasi-cohérente. En effet, soit \(E\in\Ob D(f)_{\mathrm{coh}}\); la question étant locale, on peut supposer que \(f\) est une immersion fermée, alors \(f_*(E)\) est pseudo-cohérent, donc à cohomologie quasi-cohérente, et l'on gagne du fait qu'un \(\mathcal O_X\)-Module \(M\) est quasi-cohérent si et seulement si \(f_*(M)\) l'est.

\item Quand il s'agira d'un \(\mathcal O_X\)-Module, nous utiliserons la terminologie « de présentation finie rel. à \(f\) » de préférence à « \((-1)\)-pseudo-cohérent rel. à \(f\) ». Pour qu'un \(\mathcal O_X\)-Module \(E\) soit de présentation finie rel. à \(f\), il faut et il suffit, compte tenu de (I 2.9), qu'il existe localement une immersion fermée \(i:X\to X'\), avec \(X'\) lisse sur \(Y\), telle que \(i_*(E)\) soit de présentation finie (au sens ordinaire). En particulier, dire que \(f\) est \((-1)\)-
